% FULL 49-page edition of the report (the version first submitted to GitHub, commit 39be49f).
% The short 21-page edition is SteamDataAnalysisReport.tex.
% Steam Game Recommendations — Data Analysis using Apache Hadoop and Pig
% Lab report (CSE 4346). Compile from this folder:  tectonic SteamDataAnalysisReport.tex
\documentclass[11pt,a4paper]{article}

\usepackage[a4paper,margin=2.2cm]{geometry}
\usepackage{fontspec}
\setmainfont{texgyretermes}[Extension=.otf, UprightFont=*-regular, BoldFont=*-bold,
                            ItalicFont=*-italic, BoldItalicFont=*-bolditalic]
\setmonofont{Inconsolatazi4}[Extension=.otf, UprightFont=*-Regular, BoldFont=*-Bold, Scale=0.92]
\usepackage{xcolor}
\usepackage{graphicx}
\usepackage{float}
\usepackage{booktabs}
\usepackage{tabularx}
\usepackage{array}
\usepackage{enumitem}
\usepackage{listings}
\usepackage[normalem]{ulem}
\usepackage{titlesec}
\usepackage{needspace}
\usepackage{caption}
\usepackage{fancyhdr}
\usepackage{pdfpages}
\usepackage[hidelinks]{hyperref}

\definecolor{secblue}{HTML}{2E74B5}
\definecolor{jobred}{HTML}{C00000}
\definecolor{subgreen}{HTML}{00A650}
\definecolor{kw}{RGB}{127,0,85}
\definecolor{cm}{RGB}{63,127,95}
\definecolor{st}{RGB}{42,0,255}
\definecolor{codebg}{HTML}{F7F7F5}
\definecolor{rule}{HTML}{D9D9D6}

% Black-and-white print edition: SteamDataAnalysisReport-print.tex defines \printversion
% and points \shotdir / \chartdir at the black-on-white screenshots and charts.
\providecommand{\shotdir}{../screenshots}
\providecommand{\chartdir}{../visualizations}
\ifdefined\printversion
  \definecolor{secblue}{HTML}{000000}
  \definecolor{jobred}{HTML}{000000}
  \definecolor{subgreen}{HTML}{000000}
  \definecolor{kw}{HTML}{000000}
  \definecolor{cm}{HTML}{555555}
  \definecolor{st}{HTML}{333333}
  \definecolor{codebg}{HTML}{FFFFFF}
  \definecolor{rule}{HTML}{777777}
\fi

\titleformat{\section}{\needspace{6\baselineskip}\color{secblue}\LARGE\bfseries}{\thesection.}{0.5em}{\uline}
\titleformat{\subsection}{\needspace{6\baselineskip}\color{jobred}\Large\bfseries}{\thesubsection}{0.5em}{\uline}
\titleformat{\subsubsection}{\needspace{5\baselineskip}\color{subgreen}\large\bfseries}{}{0em}{}
\titlespacing*{\subsubsection}{0pt}{1.4ex plus .4ex}{0.6ex}
\setcounter{tocdepth}{2}
\setlength{\parskip}{0.55em}
\setlength{\parindent}{0pt}
\setlist{itemsep=0.15em, topsep=0.3em}
\captionsetup{font=small, labelfont=bf, skip=4pt}

\pagestyle{fancy}
\fancyhf{}
\fancyhead[L]{\small\color{gray}Steam Data Analysis using Hadoop}
\fancyhead[R]{\small\color{gray}CSE 4346 — Big Data Analysis Laboratory}
\fancyfoot[C]{\small\thepage}
\renewcommand{\headrulewidth}{0.3pt}

\lstdefinelanguage{Pig}{
  morekeywords={LOAD,USING,AS,FILTER,BY,IS,NOT,NULL,GROUP,ALL,FOREACH,GENERATE,ORDER,DESC,
                LIMIT,FLATTEN,STORE,INTO,JOIN,COUNT,SUM,SPRINTF,PigStorage,int,chararray,double},
  sensitive=false, morecomment=[l]{--}, morestring=[b]'
}
\lstset{
  basicstyle=\ttfamily\footnotesize, keywordstyle=\color{kw}\bfseries, commentstyle=\color{cm},
  stringstyle=\color{st}, backgroundcolor=\color{codebg}, frame=single, rulecolor=\color{rule},
  framesep=4pt, xleftmargin=4pt, xrightmargin=4pt, breaklines=true, breakatwhitespace=true,
  columns=fullflexible, keepspaces=true, showstringspaces=false, tabsize=4, upquote=true,
  numbers=left, numberstyle=\tiny\color{gray}, numbersep=6pt, aboveskip=0.6em, belowskip=0.6em
}
\lstset{language=Java}

% \shot{file}{caption}: a terminal or browser screenshot, never taller than the page
\newcommand{\shot}[2]{%
  \begin{figure}[H]\centering
  \includegraphics[width=\linewidth,height=0.80\textheight,keepaspectratio]{#1}
  \caption{#2}\end{figure}}
\newcommand{\code}[1]{\texttt{#1}}
% a Pig script is short: never split it across two pages
\newcommand{\pigscript}[1]{\par\noindent\begin{minipage}{\linewidth}
  \subsubsection*{Script}\lstinputlisting[language=Pig]{#1}\end{minipage}\par}

\begin{document}

% cover page (made in Cover-Pages/, Word + PDF)

\includepdf[pages=1]{../report/cover.pdf}

\tableofcontents
\newpage

% ---------------------------------------------------------------------------
\section{Problem Statement}

Analyse the Steam game-recommendations dataset with Apache Hadoop (MapReduce) and Apache Pig,
across its game and review tables, to find which games and rating groups draw the most reviews and
the most hours of play on the Steam platform.

The work is split the way the course asks for it:

\begin{itemize}
  \item \textbf{Three MapReduce programs in Java}, each built on a different MapReduce design
        pattern — summarization, data organization (binning) and a reduce-side join.
  \item \textbf{Five Apache Pig Latin scripts} that ask ranking and grouping questions of the same
        data, including a map-side replicated join over the full 41-million-row review table.
\end{itemize}

Everything was run on a single-node Hadoop cluster, and every job result in this report was read back
out of HDFS after the job that produced it had finished.

% ---------------------------------------------------------------------------
\section{Dataset Description}

\textbf{Dataset:} Game Recommendations on Steam (Kaggle)\\
\url{https://www.kaggle.com/datasets/antonkozyriev/game-recommendations-on-steam}

The dataset is a crawl of the Steam store and its user reviews. It comes as three related CSV
tables and one JSON file of game descriptions. \code{app\_id} links a review to its game and
\code{user\_id} links it to its author. The analyses below use \code{games.csv} and
\code{recommendations.csv}; \code{users.csv} is described for completeness.

\begin{table}[H]\centering
\begin{tabular}{lrrl}
\toprule
File & Size & Rows & One row is \\
\midrule
\code{recommendations.csv} & 2.0 GB & 41,154,794 & one user review of one game \\
\code{users.csv}           & 193 MB & 14,306,064 & one Steam user \\
\code{games.csv}           & 4.9 MB & 50,872     & one game on the Steam store \\
\code{games\_metadata.json} & 18 MB & 50,872    & description and tags of a game (not used) \\
\bottomrule
\end{tabular}
\caption{Files in the dataset (row counts exclude the header line).}
\end{table}

All three CSV files are comma separated, start with a header row and contain no quoted fields.
Before any job was written I checked that every row has the expected number of columns
(13 for games, 3 for users, 8 for recommendations); none failed, so the files needed no cleaning
step beyond skipping the header line.

\subsection*{games.csv}
\begin{table}[H]\centering\small
\begin{tabularx}{\linewidth}{rlll X}
\toprule
\# & Field & Type & Example & Meaning \\
\midrule
0  & \code{app\_id}         & int     & 13500 & Steam application id (key) \\
1  & \code{title}           & string  & Prince of Persia: Warrior Within™ & game title \\
2  & \code{date\_release}   & date    & 2008-11-21 & release date \\
3–5 & \code{win, mac, linux} & boolean & true, false, false & supported platforms \\
6  & \code{rating}          & string  & Very Positive & Steam's review summary label (9 values) \\
7  & \code{positive\_ratio} & int     & 84 & \% of Steam reviews that are positive \\
8  & \code{user\_reviews}   & int     & 2199 & number of reviews on Steam \\
9  & \code{price\_final}    & float   & 9.99 & price in USD after discount \\
10 & \code{price\_original} & float   & 9.99 & price in USD before discount \\
11 & \code{discount}        & float   & 0.0 & discount in \% \\
12 & \code{steam\_deck}     & boolean & true & runs on the Steam Deck \\
\bottomrule
\end{tabularx}
\end{table}

\subsection*{users.csv}
\begin{table}[H]\centering\small
\begin{tabular}{rlll}
\toprule
\# & Field & Type & Meaning \\
\midrule
0 & \code{user\_id}  & int & anonymised user id (key) \\
1 & \code{products}  & int & number of games the user owns \\
2 & \code{reviews}   & int & number of reviews the user has written \\
\bottomrule
\end{tabular}
\end{table}

\subsection*{recommendations.csv}
\begin{table}[H]\centering\small
\begin{tabular}{rllll}
\toprule
\# & Field & Type & Example & Meaning \\
\midrule
0 & \code{app\_id}          & int     & 975370     & game reviewed $\rightarrow$ \code{games.app\_id} \\
1 & \code{helpful}          & int     & 0          & users who marked the review helpful \\
2 & \code{funny}            & int     & 0          & users who marked the review funny \\
3 & \code{date}             & date    & 2022-12-12 & review date \\
4 & \code{is\_recommended}  & boolean & true       & thumbs up or thumbs down \\
5 & \code{hours}            & float   & 36.3       & hours the reviewer had played the game \\
6 & \code{user\_id}         & int     & 51580      & reviewer $\rightarrow$ \code{users.user\_id} \\
7 & \code{review\_id}       & int     & 0          & unique review id \\
\bottomrule
\end{tabular}
\end{table}

The \code{rating} column is the label Steam shows on a store page. It takes nine values, and it is
the category field used for grouping and binning throughout this project, the same role the video
category played in the YouTube lab.

% ---------------------------------------------------------------------------
\section{Summary of Data Analysis}

The following MapReduce design patterns are implemented:

\begin{enumerate}
  \item \textbf{Summarization} — numerical summarization with a custom \code{Writable} and a combiner.
  \item \textbf{Data Organization} — binning with \code{MultipleOutputs}.
  \item \textbf{Join Patterns} — reduce-side join with \code{MultipleInputs}.
\end{enumerate}

The following MapReduce analyses are performed on the dataset:

\begin{enumerate}
  \item \textbf{Review summary per game} — number of reviews, recommended share, total, average and
        maximum hours played, and helpful votes for every game (\emph{summarization}).
  \item \textbf{Binning games by rating} — every game written to one of nine output files, one per
        Steam rating label (\emph{data organization}).
  \item \textbf{Games joined with their reviews} — each game's title and rating put next to its review
        count and total hours played (\emph{reduce-side join}).
\end{enumerate}

The following Pig analyses are performed on the dataset:

\begin{enumerate}
  \item Top 5 rating groups by number of games
  \item Top 10 most reviewed games
  \item Top 10 most reviewed games in each rating group
  \item Top 10 games by total hours played
  \item Top 10 games by total hours played in each rating group
\end{enumerate}

% ---------------------------------------------------------------------------
\section{Environment and Loading the Data}

\begin{table}[H]\centering
\begin{tabular}{ll}
\toprule
Component & Version / setting \\
\midrule
Operating system & Pop!\_OS Linux, 8 CPU cores, 16 GB RAM \\
Apache Hadoop & 3.4.1, pseudo-distributed (single node), \code{fs.defaultFS = hdfs://localhost:9000} \\
HDFS & replication 1, block size 128 MB \\
Apache Pig & 0.18.0, run in \code{mapreduce} mode on YARN \\
Java & OpenJDK 11.0.31 \\
YARN memory & NodeManager 5120 MB; map 768 MB, reduce 1024 MB, AM 768 MB containers \\
\bottomrule
\end{tabular}
\caption{Software used.}
\end{table}

The laptop had only about 2 GB of RAM free, and with Hadoop's default settings YARN would try to run
eight 1 GB task containers at once. The NodeManager was therefore capped at 5 GB and the task
containers made smaller in \code{yarn-site.xml} and \code{mapred-site.xml}, so the jobs run a few
tasks at a time instead of pushing the machine into swap.

\subsubsection*{Starting the cluster}
The cluster is started daemon by daemon, followed by the JobHistory server. \code{jps} shows all six
Hadoop daemons running:

\shot{\shotdir/01_cluster_jps.png}{All six Hadoop daemons running (\code{jps}).}
\shot{\shotdir/02_versions.png}{Hadoop, Pig and Java versions.}

\subsubsection*{Loading the dataset into HDFS}
The three CSV files are copied into the \code{/steam} directory in HDFS. The 2 GB review file is
split into 16 blocks of 128 MB, and \code{hdfs fsck} reports the directory as healthy.

\shot{\shotdir/03_hdfs_put.png}{Copying the dataset into HDFS (11.5 seconds).}
\shot{\shotdir/04_hdfs_ls.png}{The dataset in HDFS: 1.9 GB of reviews stored as 16 blocks.}

\subsubsection*{Building the MapReduce jobs}
Each job is compiled against the local Hadoop classpath and packed into its own jar by
\code{scripts/build.sh}:

\lstinputlisting[language=bash]{../scripts/build.sh}
\shot{\shotdir/05_build_jars.png}{The three jars built.}

% ---------------------------------------------------------------------------
\section{Map Reduce Analysis}

\subsection{Review Summary per Game}
\textbf{Returns, for every game, how many reviews it received, how many recommend it, how long its
reviewers played it and how many helpful votes its reviews collected} (\code{GameReviewSummary}).

This is the \textbf{numerical summarization} pattern. The mapper reads one review at a time and emits
the game's \code{app\_id} with a small tuple holding a count of one, a 1 or 0 for ``recommended'',
the hours played and the helpful votes. Every field of the tuple is additive — a count, a sum or a
maximum — so partial results can be merged in any order. That is what makes a \textbf{combiner}
valid here: each mapper merges its own tuples before the shuffle, and the reducer finishes the job
with exactly the same code.

This analysis consists of:
\begin{enumerate}
  \item \textbf{Driver Class} — sets up the job, and registers the reducer as the combiner as well.
  \item \textbf{Mapper Class} — emits \code{app\_id} of type \code{Text} and a tuple of custom type
        \code{ReviewStatsTuple}.
  \item \textbf{Reducer Class} — merges all tuples of one game and emits \code{app\_id} with the
        final \code{ReviewStatsTuple}; also used as the combiner.
  \item \textbf{Tuple Class} — custom \code{Writable} holding reviews, recommended, total hours,
        maximum hours and helpful votes.
\end{enumerate}

\subsubsection*{Driver Class}
\lstinputlisting{../GameReviewSummary/src/main/java/steam/summary/GameReviewSummary.java}
\subsubsection*{Mapper Class}
\lstinputlisting{../GameReviewSummary/src/main/java/steam/summary/ReviewMapper.java}
\subsubsection*{Reducer Class (also the Combiner)}
\lstinputlisting{../GameReviewSummary/src/main/java/steam/summary/ReviewReducer.java}
\subsubsection*{Tuple Class}
\lstinputlisting{../GameReviewSummary/src/main/java/steam/summary/ReviewStatsTuple.java}

\subsubsection*{Execution}
\shot{\shotdir/MR1_execution.png}{Review summary job running over the 2 GB review file.}
\shot{\shotdir/MR1_counters.png}{Job counters of the review summary job.}

The job read all 41,154,794 reviews in 15 input splits and finished in 2 minutes 15 seconds. The
counters show the combiner at work: the mappers produced 41,154,794 records, but after combining
only 90,766 of them reached the reducer — roughly 450 times less data through the shuffle. The
reducer wrote one line for each of the \textbf{37,610 games} that have at least one review.

\subsubsection*{Output}
\shot{\shotdir/MR1_output.png}{The ten games with the most reviews. Columns: app\_id, reviews,
recommended, positive \%, total hours, average hours, maximum hours, helpful votes.}

\begin{table}[H]\centering\small
\begin{tabular}{rlrrrr}
\toprule
app\_id & Game & Reviews & Positive \% & Avg.\ hours & Helpful votes \\
\midrule
440     & Team Fortress 2                  & 319,492 & 92.3 & 318.7 & 521,978 \\
252490  & Rust                             & 270,684 & 83.6 & 343.4 & 771,227 \\
1091500 & Cyberpunk 2077                   & 226,414 & 74.5 & 102.2 & 1,189,869 \\
730     & Counter-Strike: Global Offensive & 219,737 & 84.8 & 429.0 & 1,910,071 \\
570     & Dota 2                           & 216,914 & 79.2 & 429.1 & 842,346 \\
444090  & Paladins®                        & 204,176 & 84.0 & 139.8 & 188,983 \\
292030  & The Witcher® 3: Wild Hunt        & 204,166 & 94.6 & 154.4 & 813,793 \\
377160  & Fallout 4                        & 196,789 & 81.6 & 204.0 & 899,972 \\
374320  & DARK SOULS™ III                  & 190,157 & 92.9 & 170.7 & 529,674 \\
431960  & Wallpaper Engine                 & 190,129 & 98.2 & 102.8 & 802,753 \\
\bottomrule
\end{tabular}
\caption{The same ten games with their titles looked up from \code{games.csv}.}
\end{table}

The maximum-hours column reads 999.9 for almost every popular game. That is not a bug in the job: the
\code{hours} field in the source data never goes above 1,000, so for any game with a large player
base at least one reviewer sits at the cap.

% ---------------------------------------------------------------------------
\subsection{Binning Games by Rating}
\textbf{Performs binning on the Steam rating label, writing every game into one output file per
bin.} The program is \code{BinningByRating}.

This is the \textbf{binning} pattern from the data organization family. The job is map-only: each
mapper decides which bin a record belongs to and writes it straight to that bin's file through
\code{MultipleOutputs}, so there is no shuffle, no sort and no reducer. The bin name comes from the
\code{rating} column with the spaces removed (``Very Positive'' becomes \code{VeryPositive-m-00000}),
and a counter per bin records how many games landed in it. \code{LazyOutputFormat} stops Hadoop from
also writing an empty default \code{part-m-00000} file.

This analysis consists of:
\begin{enumerate}
  \item \textbf{Driver Class} — sets zero reducers and registers the \code{bins} named output.
  \item \textbf{Mapper Class} — uses \code{MultipleOutputs} to write each game to the file of its rating
        bin, and counts the games per bin.
\end{enumerate}

\subsubsection*{Driver Class}
\lstinputlisting{../BinningByRating/src/main/java/steam/binning/BinningByRating.java}
\subsubsection*{Mapper Class}
\lstinputlisting{../BinningByRating/src/main/java/steam/binning/BinningMapper.java}

\subsubsection*{Execution}
\shot{\shotdir/MR2_execution.png}{Binning job running on \code{games.csv}.}
\shot{\shotdir/MR2_counters.png}{Job counters, including the nine \code{Bins} counters.}

One map task read the 50,873 lines of \code{games.csv} (50,872 games and the header). The
\code{Map output records} counter is 0 because \code{MultipleOutputs} writes to the bin files
directly instead of going through the normal map output. The \code{Bins} counters are the real
result: the nine of them add up to 50,872, so every game was placed in exactly one bin.

\subsubsection*{Output}
\shot{\shotdir/MR2_output.png}{The nine bin files in HDFS, and the smallest bin
(Overwhelmingly Negative, 14 games) printed in full.}

\begin{table}[H]\centering\small
\begin{tabular}{lrr}
\toprule
Bin (rating) & Games & Share \\
\midrule
Positive                & 13,502 & 26.5\% \\
Very Positive           & 13,139 & 25.8\% \\
Mixed                   & 12,157 & 23.9\% \\
Mostly Positive         &  8,738 & 17.2\% \\
Mostly Negative         &  1,849 &  3.6\% \\
Overwhelmingly Positive &  1,110 &  2.2\% \\
Negative                &    303 &  0.6\% \\
Very Negative           &     60 &  0.1\% \\
Overwhelmingly Negative &     14 &  0.03\% \\
\midrule
Total                   & 50,872 & 100\% \\
\bottomrule
\end{tabular}
\caption{Games per bin, from the job's \code{Bins} counters.}
\end{table}

% ---------------------------------------------------------------------------
\subsection{Games Joined with their Reviews}
\textbf{Returns each game's title and rating together with how many people reviewed it and how many
hours they played it} (\code{GameReviewJoin}).

This is a \textbf{reduce-side join}. The two inputs have different layouts, so each gets its own
mapper through \code{MultipleInputs}. Both mappers use \code{app\_id} as the key and tag their values —
\code{G} for a game record, \code{R} for review data — so that the reducer, which receives
everything for one \code{app\_id} together, can tell the two sides apart and put them back together.

Sending all 41 million reviews through the shuffle one by one would be slow, so the review mapper
uses \textbf{in-mapper combining}: it keeps a running total per game in a hash map and emits a single
\code{R} record per game when its input split is done. The join is an inner join — games nobody
reviewed, and reviews whose game is missing from \code{games.csv}, are dropped and counted.

This analysis consists of:
\begin{enumerate}
  \item \textbf{Driver Class} — registers the two inputs, each with its own mapper, and one reducer.
  \item \textbf{Game Mapper} — emits \code{app\_id} and \code{"G, title, rating"}.
  \item \textbf{Recommendation Mapper} — adds up reviews, recommendations and hours per game in memory
        and emits \code{app\_id} and \code{"R, reviews, recommended, hours"}.
  \item \textbf{Join Reducer} — matches the \code{G} record with the \code{R} records of the same game and
        emits the joined line.
\end{enumerate}

\subsubsection*{Driver Class}
\lstinputlisting{../GameReviewJoin/src/main/java/steam/join/GameReviewJoin.java}
\subsubsection*{Game Mapper}
\lstinputlisting{../GameReviewJoin/src/main/java/steam/join/GameMapper.java}
\subsubsection*{Recommendation Mapper}
\lstinputlisting{../GameReviewJoin/src/main/java/steam/join/RecommendationMapper.java}
\subsubsection*{Join Reducer}
\lstinputlisting{../GameReviewJoin/src/main/java/steam/join/JoinReducer.java}

\subsubsection*{Execution}
\shot{\shotdir/MR3_execution.png}{Join job running over both inputs.}
\shot{\shotdir/MR3_counters.png}{Job counters of the join, including the \code{Join} counters.}

The job read 16 input splits — 15 from the review file and one from \code{games.csv} — and finished in
1 minute 15 seconds, faster than the summary job even though it reads the same 2 GB. The reason is
in the counters: in-mapper combining cut the map output from 41 million records to 141,638, so the
shuffle had almost nothing to move. The \code{Join} counters show 37,610 games matched and 13,262 games
with no reviews at all; no review pointed to a game missing from \code{games.csv}.

\subsubsection*{Output}
\shot{\shotdir/MR3_output.png}{The ten games with the most hours played. Columns: app\_id,
title, rating, reviews, positive \%, total hours, average hours per reviewer.}

\begin{table}[H]\centering\small
\begin{tabular}{llrrr}
\toprule
Game & Rating & Reviews & Total hours & Avg.\ hours \\
\midrule
Team Fortress 2                   & Very Positive & 319,492 & 101,809,977.6 & 318.7 \\
Counter-Strike: Global Offensive  & Very Positive & 219,737 &  94,260,307.3 & 429.0 \\
Dota 2                            & Very Positive & 216,914 &  93,072,785.4 & 429.1 \\
Rust                              & Very Positive & 270,684 &  92,944,004.6 & 343.4 \\
Tom Clancy's Rainbow Six® Siege   & Very Positive & 189,603 &  68,445,823.3 & 361.0 \\
Destiny 2                         & Very Positive & 185,527 &  62,504,172.8 & 336.9 \\
Hearts of Iron IV                 & Very Positive & 143,187 &  53,086,631.1 & 370.8 \\
ARK: Survival Evolved             & Very Positive & 159,736 &  46,720,359.4 & 292.5 \\
Warframe                          & Very Positive & 127,812 &  46,450,175.3 & 363.4 \\
Grand Theft Auto V                & Very Positive & 163,909 &  46,170,394.5 & 281.7 \\
\bottomrule
\end{tabular}
\caption{Top 10 games by total hours played, from the joined output.}
\end{table}

% ---------------------------------------------------------------------------
\section{Pig Analysis}

All five scripts read the dataset straight from HDFS with \code{PigStorage(',')} and drop the header row
with \code{FILTER \dots\ BY app\_id IS NOT NULL} — the header's \code{app\_id} is the text
``app\_id'', which fails the cast to \code{int} and becomes null. Results are stored tab separated
under \code{/steam/pig/}; a tab is used instead of a pipe because 104 game titles contain a
\code{|} character.

Every ranking sorts \emph{inside} a nested \code{FOREACH} block (\code{ORDER} then \code{LIMIT}) rather
than with a top-level \code{ORDER BY}. For the ``in each rating group'' queries this is what the
question needs — a separate top 10 per group — and for the overall rankings it avoids the
sampling job that a top-level \code{ORDER BY} starts.

Each run ends with a Pig warning, \code{FIELD\_DISCARDED\_TYPE\_CONVERSION\_FAILED}. It counts the
fields of the header rows that could not be cast to a number, and it is exactly how the header is
filtered out, so it is expected rather than a data problem.

\subsection{Top 5 Rating Groups}
\textbf{Returns the five Steam rating labels that the most games carry.}

The games are grouped by \code{rating} and counted; the nine counts are then gathered into a single
group (\code{GROUP \dots\ ALL}) so they can be sorted together and cut to five. Pig compiles the script
into two MapReduce jobs — the count, which runs with a combiner, and the ranking.

\pigscript{../PigAnalysis/scripts/top5Ratings.pig}
\subsubsection*{Execution}
\shot{\shotdir/Pig1_execution.png}{\code{top5Ratings.pig} starting in \code{mapreduce} mode.}
\shot{\shotdir/Pig1_result.png}{End of the run: \code{Success!}, the two jobs, 50,873 records read and 5 stored.}
\subsubsection*{Output}
\shot{\shotdir/Pig1_output.png}{Top 5 rating groups.}
The five counts are identical to the first five bins of the MapReduce binning job (Section 5.2).
\subsubsection*{Visualization}
\begin{figure}[H]\centering\includegraphics[width=0.85\linewidth]{\chartdir/pig1_top5_ratings.png}
\caption{Top 5 rating groups by number of games.}\end{figure}

\subsection{Top 10 Most Reviewed Games}
\textbf{Returns the ten games with the most reviews on Steam} (the \code{user\_reviews} column).

\pigscript{../PigAnalysis/scripts/top10Reviewed.pig}
\subsubsection*{Execution}
\shot{\shotdir/Pig2_execution.png}{\code{top10Reviewed.pig} starting.}
\shot{\shotdir/Pig2_result.png}{End of the run: one MapReduce job, 10 records stored in 48 seconds.}
\subsubsection*{Output}
\shot{\shotdir/Pig2_output.png}{Top 10 most reviewed games. Columns: app\_id, title, rating,
positive ratio, user reviews.}

Counter-Strike: Global Offensive has 7,494,460 reviews on Steam — more than three times the next game,
PUBG: BATTLEGROUNDS (2,217,226). This column is Steam's own total for the whole store, which is
why its order differs from MapReduce job 1: that job counts only the reviews present in
\code{recommendations.csv}, where Team Fortress 2 comes first.

\subsubsection*{Visualization}
\begin{figure}[H]\centering\includegraphics[width=0.85\linewidth]{\chartdir/pig2_top10_reviewed.png}
\caption{Top 10 most reviewed games.}\end{figure}

\subsection{Top 10 Most Reviewed Games in each Rating Group}
\textbf{Returns the ten most reviewed games inside each of the nine rating groups.}

This is the nested-\code{FOREACH} pattern: after \code{GROUP BY rating}, the block runs once per group,
sorts that group's bag by \code{user\_reviews}, keeps ten and \code{FLATTEN}s them back into rows. The
whole script runs as a single MapReduce job.

\pigscript{../PigAnalysis/scripts/top10ReviewedByRating.pig}
\subsubsection*{Execution}
\shot{\shotdir/Pig3_execution.png}{\code{top10ReviewedByRating.pig} starting.}
\shot{\shotdir/Pig3_result.png}{End of the run: 90 records stored — nine groups of ten.}
\subsubsection*{Output}
\shot{\shotdir/Pig3_output.png}{Top 10 most reviewed games in every rating group.}

Two things stand out. The ``Positive'' and ``Negative'' groups never go above 49 reviews: Steam only
gives those plain labels to games with 10–49 reviews, and a game that collects more moves up to
``Very'' or ``Overwhelmingly''. And the most reviewed game in the Overwhelmingly Negative group is
Overwatch® 2 with 181,198 reviews — a very large game can sit in the worst group.

\subsubsection*{Visualization}
\begin{figure}[H]\centering\includegraphics[width=\linewidth]{\chartdir/pig3_top10_reviewed_by_rating.png}
\caption{Top 10 most reviewed games in each rating group. Each panel has its own scale.}\end{figure}

\subsection{Top 10 Games by Total Hours Played}
\textbf{Returns the ten games whose reviewers played them for the most hours in total.}

This script reads the full 2 GB review table. It adds up the hours per game with
\code{GROUP BY app\_id} and \code{SUM}, then brings in each game's title and rating with a
\textbf{replicated join}: \code{games.csv} is small enough to be copied to every mapper, so the join
happens on the map side without a second shuffle. The hours stay numeric for the sort and are only
formatted with \code{SPRINTF} for the output.

Pig turns the script into a chain of three MapReduce jobs (visible in the Job Stats of the
screenshot below): a map-only job that prepares \code{games.csv} for the replicated join, one job with
15 maps and 3 reducers that groups the 41 million reviews and performs the join, and a final job
that ranks the result.

\pigscript{../PigAnalysis/scripts/top10Played.pig}
\subsubsection*{Execution}
\shot{\shotdir/Pig4_execution.png}{\code{top10Played.pig} starting.}
\shot{\shotdir/Pig4_result.png}{End of the run: 41,154,795 review lines and 50,873 game lines
read, 10 records stored, in 4 minutes 35 seconds.}
\subsubsection*{Output}
\shot{\shotdir/Pig4_output.png}{Top 10 games by total hours played. Columns: app\_id, title,
rating, reviews, total hours.}

The list is identical to the top ten of the MapReduce join (Section 5.3) — same games, same order,
same review counts and the same hours to one decimal place. The two tools took very different routes
to it: the hand-written join finished in 1 minute 15 seconds, the Pig script in 4 minutes 35 seconds,
mostly because Pig runs three jobs one after another while the Java version does everything in one
job and shrinks its shuffle with in-mapper combining.

\subsubsection*{Visualization}
\begin{figure}[H]\centering\includegraphics[width=0.85\linewidth]{\chartdir/pig4_top10_played.png}
\caption{Top 10 games by total hours played by reviewers.}\end{figure}

\subsection{Top 10 Games by Total Hours Played in each Rating Group}
\textbf{Returns the ten most played games inside each rating group.}

Same first half as the previous script — sum the hours per game and replicated-join the titles and
ratings — followed by the nested sort of Section 6.3, grouped by rating.

\pigscript{../PigAnalysis/scripts/top10PlayedByRating.pig}
\subsubsection*{Execution}
\shot{\shotdir/Pig5_execution.png}{\code{top10PlayedByRating.pig} starting.}
\shot{\shotdir/Pig5_result.png}{End of the run: 87 records stored, in 4 minutes 3 seconds.}
\subsubsection*{Output}
\shot{\shotdir/Pig5_output.png}{Top 10 most played games in every rating group.}

The script stored 87 rows, not 90. The Overwhelmingly Negative group has 14 games, but only 7 of them
have any review in \code{recommendations.csv}, so that group can only return seven.

\subsubsection*{Visualization}
\begin{figure}[H]\centering\includegraphics[width=\linewidth]{\chartdir/pig5_top10_played_by_rating.png}
\caption{Top 10 games by total hours played in each rating group. Each panel has its own scale.}\end{figure}

% ---------------------------------------------------------------------------
\section{Monitoring the Jobs}

Hadoop's web interfaces were kept open while the jobs ran.

\shot{\shotdir/web_yarn_applications.png}{YARN ResourceManager (\code{localhost:8088}): all 19
applications finished with status SUCCEEDED.}

The 19 applications are the three MapReduce jobs and the MapReduce jobs Pig generated for the five
scripts. Scripts 4 and 5 appear twice: they were run once, then run again after changing only how the
hours are printed (from Pig's default scientific notation to one decimal place). The second runs are
the ones in this report.

\shot{\shotdir/web_jobhistory.png}{MapReduce JobHistory server (\code{localhost:19888}) with the
map and reduce task counts of every job.}
\shot{\shotdir/web_namenode_overview.png}{NameNode overview (\code{localhost:9870}).}
\shot{\shotdir/web_hdfs_steam.png}{The \code{/steam} directory in the HDFS browser: the three input
files and the output directories of every job.}

% ---------------------------------------------------------------------------
\section{Verification}

Several results can be checked against each other or against the raw file, because different jobs
reach the same number by different routes. All checks passed.

\begin{table}[H]\centering\small
\begin{tabularx}{\linewidth}{X r r c}
\toprule
Check & Expected & Produced & \\
\midrule
Reviews read by the summary job = rows of \code{recommendations.csv} & 41,154,794 & 41,154,794 & pass \\
Reviews summed over the summary output = over the join output & 41,154,794 & 41,154,794 & pass \\
Games in the summary output = games matched by the join & 37,610 & 37,610 & pass \\
Joined games + games with no reviews = rows of \code{games.csv} & 50,872 & 37,610 + 13,262 & pass \\
Nine bins of the binning job add up to all games & 50,872 & 50,872 & pass \\
Bin sizes = counts computed directly from \code{games.csv} & 9 values & all 9 equal & pass \\
Pig top 5 rating groups = the five largest bins & 5 values & all 5 equal & pass \\
Pig top 10 by hours = MapReduce join top 10 (ids, reviews, hours) & 10 rows & identical & pass \\
Pig-5 rows = 8 groups $\times$ 10 + 7 reviewed Overwhelmingly Negative games & 87 & 87 & pass \\
\bottomrule
\end{tabularx}
\caption{Cross-checks between jobs.}
\end{table}

The bin sizes were also counted outside Hadoop, straight from the CSV, before any job was run, and
the Pig-versus-MapReduce comparison of the top ten was done with \code{diff} on the two outputs.

% ---------------------------------------------------------------------------
\section{What the Data Says}

\textbf{Steam reviews are mostly positive.} 85.8\% of the 41 million reviews recommend the game, and
71.7\% of all games carry one of the four positive labels. Only 2,226 games (4.4\%) carry a negative
one, and just 14 are Overwhelmingly Negative.

\textbf{Attention is concentrated in a few games.} Of the 50,872 games, 37,610 appear in the review
data at all, and half of those have 39 reviews or fewer. The ten most reviewed games — 0.03\% of the
reviewed games — hold 5.4\% of all reviews.

\textbf{The most played games are long-running multiplayer titles.} Team Fortress 2, Counter-Strike:
Global Offensive, Dota 2 and Rust each passed 90 million hours of play among their reviewers, and the
average Counter-Strike or Dota 2 reviewer had played about 429 hours. Across all reviews the players
had logged 4.14 billion hours.

\textbf{Popular does not mean well rated.} PUBG: BATTLEGROUNDS is the second most reviewed game on
Steam but is rated Mixed, with 57\% positive reviews; Overwatch® 2 has 181,198 reviews and is
Overwhelmingly Negative. At the other end, Wallpaper Engine is 98.2\% positive in the review data.

\textbf{A rating label is partly a volume label.} The Positive and Negative groups top out at 49
reviews, because Steam reserves the ``Very'' and ``Overwhelmingly'' labels for games with more reviews.
Grouping by \code{rating} therefore mixes two things — how good the reviews are and how many there are
— which is worth keeping in mind before comparing the groups directly.

% ---------------------------------------------------------------------------
\section{Conclusion}

A 2.2 GB, three-table Steam dataset was loaded into HDFS and analysed on a single-node Hadoop 3.4.1
cluster with three Java MapReduce programs and five Pig Latin scripts, all of which completed
successfully.

The three MapReduce programs each show a different design pattern. The summarization job shows how a
custom \code{Writable} and a combiner shrink 41 million map outputs to under a hundred thousand before
the shuffle. The binning job shows a map-only job writing many outputs at once. The reduce-side join
shows two differently shaped inputs meeting in one reducer, and how in-mapper combining makes the
join over 41 million rows finish in just over a minute.

Pig answered the ranking questions in 15 to 33 lines per script, including a replicated join over the
full review table, but took roughly three to four times as long as the hand-tuned Java join for the
same answer. The two tools agreed on every result they share, which is the strongest evidence that
both are correct.

% ---------------------------------------------------------------------------
\section*{References}
\addcontentsline{toc}{section}{References}
\begin{enumerate}[label={[\arabic*]}]
  \item A.\ Kozyriev, \emph{Game Recommendations on Steam}, Kaggle dataset.
        \url{https://www.kaggle.com/datasets/antonkozyriev/game-recommendations-on-steam}
  \item S.\ Prabhu, \emph{YouTube Data Analysis using Hadoop} (project structure followed).
        \url{https://github.com/SiddhiPrabhu1995/YouTube-Data-Analysis-using-Hadoop}
  \item D.\ Miner and A.\ Shook, \emph{MapReduce Design Patterns}, O'Reilly, 2012.
  \item Apache Hadoop 3.4.1 documentation. \url{https://hadoop.apache.org/docs/r3.4.1/}
  \item Apache Pig 0.18.0 documentation. \url{https://pig.apache.org/docs/latest/}
\end{enumerate}

\end{document}
